\providecommand{\PushA}{\frac{5422587}{6103515625}}
\providecommand{\PushExponent}{\frac{5422587}{433349609375}}
\providecommand{\PushK}{890666}
\providecommand{\PushTheta}{0.99999999997327}
\section{Baseline closing: \texorpdfstring{$\Theta<1$}{Theta<1} with \texorpdfstring{$\theta=1$}{theta=1}}\label{sec:closing}

\paragraph{What this section proves and how to read it.}
This section proves the baseline \cref{thm:baseline} only. Recall
$\Theta=(63/64)(L_0+\mathcal E)+J+W+(1-m)/64$ with
$\mathcal E=B(e^S\cosh A-1)-D$ (\cref{thm:comparison} and
\cref{subsec:intro:inequality}). Every quantity is defined at
$a=\PushA$ as a sum of middle prices, upper-band prices converted
from their evaluation exponents, or an explicit function thereof.
Outward-interval evaluation gives $\Theta<\PushTheta$; Theorem~C
with $\theta=1$ then gives $\alpha^{71}\ll\Lsc^{-a}$, i.e.
$\eps=a/71$, the $71$st-root step being
$\alpha\le C\Lsc^{-a/71}$ with an absolute constant $C$. The
margin $1-\Theta$ is about
$2.7\cdot10^{-11}$, so numbers are stated exactly as evaluated; no bound may be replaced by a larger bound for the same quantity without re-evaluating $\Theta$. The
words \emph{robust} and \emph{inflation} below mean: errors come
only from the two allowed remainders (\cref{def:robust}), and
prices are increased to the theorem exponent $a$ from smaller
evaluation exponents (repricing). The improved theorem redoes this
evaluation at $\vartheta>1$ in \cref{sec:four-routes-closing}.

\subsection{Definitions of \texorpdfstring{$L_0,S,A,D,\mathcal E,T,R_I,J,W,m$}{L\_0,S,A,D,E,T,R\_I,J,W,m} (baseline)}\label{subsec:closing-quantities}

This subsection defines the ten numbers entering $\Theta$. Two
conversions are applied first. The middle prices were computed
(certified) at a smaller exponent $a_0=\PushAnchor<a$, the
\emph{anchor}, and are converted to the exponent $a$ of the
theorem (``repriced'') by \cref{cor:first-stage-repricing} and
\cref{lem:record:repricing}, using the block bound
$\sum_{p\in\Blk}\log p\le L=570$. The upper prices have their own
anchor $a_{\rm u}=883/10^6$ and are repriced term by term in each
band (\cref{prop:record:upper}). All rounding is upward: we write
$\operatorname{up}(x)=10^{-30}\lceil10^{30}x\rceil$, the next
multiple of $10^{-30}$ above $x$. A certified quantity $v$ is
stored in the data file (the ``native source'' of
\cref{app:certificates}) as an interval enclosure with an exact
rational upper end $v^{\rm ex}$ and a decimal upper end
$v^{\rm hi}\ge v^{\rm ex}$; we put
$\operatorname{hi}(v)=\operatorname{up}(v^{\rm hi})$, and it is
these rounded upper ends that enter every formula below. For the middle piece $i$
($1\le i\le33$), with left endpoint $X_i$ of its prime interval (\cref{sec:middle-reference}) and $q_{X_i}=1/X_i$, and for $d\in\{2,3,4\}$,
$I_{d,i}$ is the certified degree-$d$ charge of the piece at $a_0$, the bound
of \cref{lem:record:block-charge} for
$\sum_{\Blk}q_{\Blk}^{2a_0}b_{\Blk}^{-d/2}$ ($b_\Blk$ the least
prime of $\Blk$) over its blocks
(\cref{cert:record:middle-price}), and $\bar I_{d,i}=\operatorname{hi}(I_{d,i})$.
A local coefficient evaluated at $X_i$ enters a price through
$\operatorname{up}$ of the upper end of its outward enclosure;
we write this rounded majorant with the same letter as the
coefficient. Thus every letter below denotes the rounded upper
value, and every inequality is in the safe (upper) direction.

\begin{definition}[The closing quantities]\label{def:closing-quantities}
\leavevmode
\begin{enumerate}[label=(\arabic*)]
\item \emph{Repricing.} $a_0=\PushAnchor$, $L=570$ and $t=e^{2(a-a_0)L}$;
 $C=e^{2aL}$; the piece caps $C_i=e^{2aL_i}$ of \cref{lem:piece-caps}, and
 $C_X=\prod_{\Blk\in X}C_{i(\Blk)}$ for a set $X$ of middle blocks, $i(\Blk)$
 the piece of $\Blk$.
\item \emph{First stage and bracket.} Let $U_0'$, $P'$, $D'$ and $V'$ be the
 recorded upper ends of the totals at $a_0$ of \cref{lem:record:middle-prices}
 over the $33$ pieces (\cref{cert:record:middle-price}):
 $U_0'\le0.8469081935391429860$, $P'\le1.6936504653996419205$,
 $D'\le4.741704122562283\cdot10^{-5}$ and $V'\le3.3872533701481774716$. Put
 \[
  U_0=tU_0',\qquad P=tP',\qquad d=\tfrac12tD',\qquad v=\tfrac12tV',
 \]
 where $D'$ and $V'$ are the stored totals at $a_0$ of the sharp prices \eqref{eq:record:sharp} and total-variation prices summed over all middle atoms and both first-stage record differences (the ``tags'' of \cref{def:record:frame}(e)), and the two tags carry equal halves of these totals by \cref{lem:record:middle-prices}, so that $d=\tfrac12tD'$ and $v=\tfrac12tV'$ are the repriced per-tag totals $d_j,v_j$ of \cref{def:cp2}, and
 \[
  R_{\rm err}=e^{U_0}\tagser_4^+(d,v),\qquad L_0=e^{P}\tagser_4^+(d,v),\qquad
  B=e^{U_0}+R_{\rm err},
 \]
 with the majorant $\tagser_4^+$ of \cref{cor:two-cut-series}. By
 \cref{prop:closing-data}(a) these numbers bound the totals of the same names
 of (CP2) and (CP3), and the closing uses them in place of those totals
 (\cref{rem:cp:monotone}); \cref{sec:variance} writes $B_*$ and $R_*$ for them.
 (The letter $D$ is reserved for the discount of (10).)
\item \emph{Middle prices.} For the piece $i$ the \emph{rows} at $a_0$ are the summands of \cref{tab:data:prices} for that piece, as stored in the data, before the normalisation by $Z_i^{-1}$ of \cref{sec:four-routes-closing}. The fold and matching rows agree with
 \eqref{eq:middle-piece-prices}; the two generic rows $S_{{\rm oth},i}$ and $A_{{\rm oth},i}$ use the larger caps stored in the data (\cref{tab:data:prices}), at least those of \eqref{eq:middle-piece-prices}, which also bound the two-layer
 constant and merge repairs.
 Explicitly, for its two
 fold records the raw-arm ($y^{\rm arm}$), leading, remainder, full skew (leading plus remainder) and symmetric prices
 $A^{\rm arm}_{F,i}$, $A^{\rm lead}_{F,i}$, $A^{\rm rem}_{F,i}$,
 $A_{F,i}=A^{\rm lead}_{F,i}+A^{\rm rem}_{F,i}$ and $S_{F,i}$; the same prices
 $A^{\rm arm}_{E,i}$, $A^{\rm lead}_{E,i}$, $A^{\rm rem}_{E,i}$, $A_{E,i}$, $S_{E,i}$
 of its two matching records; and the symmetric and skew prices $S_{{\rm oth},i}$,
 $A_{{\rm oth},i}$ of its other records (merges, deletions and the remaining
 changed records). They are formed from the completion weights
 $w_{{\rm D},i}$, $w_{{\rm S},i}$ (\cref{rem:middle-weights}), the charges
 $\bar I_{d,i}$, the endpoint coefficients $\ell_\tlaw(X_i)$, $a_F(X_i)$, $s_F(X_i)$
 of the signed fold records (\cref{prop:signed-bounds}), $a_E(X_i)$, $s_E(X_i)$,
 $r_E(X_i)$ of the matching records
 (\cref{lem:record:exact-marginal-e,lem:oct24:matching-star}) and the source coefficients for the generic rows $S_{{\rm oth},i},A_{{\rm oth},i}$, and rounded per piece on
 $10^{-30}\Z$; for instance
 $A^{\rm lead}_{F,i}=\operatorname{up}(2w_{{\rm D},i}\ell_\tlaw(X_i)\bar I_{2,i})$ and
 $A_{E,i}=2\operatorname{up}(w_{{\rm S},i}a_E(X_i)\bar I_{2,i})$. The totals at
 $a$ are $t$ times the sums over the pieces:
 $A^{\rm arm}_F=t\sum_iA^{\rm arm}_{F,i}$, and likewise $A^{\rm lead}_F$,
 $A^{\rm rem}_F$, $A_F$, $S_F$, $A^{\rm arm}_E$, $A^{\rm lead}_E$,
 $A^{\rm rem}_E$, $A_E$, $S_E$, $S_{\rm oth}$ and $A_{\rm oth}$; and
 $S_{\rm mid}=S_F+S_E+S_{\rm oth}$, $A_{\rm mid}=A_F+A_E+A_{\rm oth}$.
\item \emph{Symmetric and skew totals.}
 \[
  S=S_{\rm mid}+S_{\rm u}(a)+10^{-12},\qquad A=A_{\rm mid}+A_{\rm u}(a)+10^{-12},
 \]
 where $S_{\rm u}(a)$, $A_{\rm u}(a)$ are the symmetric and skew ledgers of the
 $310$ upper bands of \cref{prop:record:upper}, certified at $a_{\rm u}$ and
 repriced term by term, each term with its own logarithmic product bound, plus
 the tail sums of \cref{lem:tail-prices}(iv).
\item \emph{Matching exponent.} $T_K=t\bigl(T^{\rm aux\prime}_K+T^{\rm inc\prime}_K\bigr)$
 is the price of the chosen side diagonals of \cref{prop:signed-sides}(iii):
 $T^{\rm aux\prime}_K=\sum_i\operatorname{up}\bigl((k^{\rm aux}_{{\rm L},i}
 +k^{\rm aux}_{{\rm R},i})\bar I_{2,i}\bigr)$ is the price bound for the excess coefficients
 $k^{\rm aux}_{s,i}=w_{{\rm D},i}d^{\rm aux}_s(X_i)-v_{{\rm D},i}$ of the auxiliary
 sides (orientation weights, local cap $L_{\rm loc}=585$), and $T^{\rm inc\prime}_K$
 the increments $\delta^+_{s,i}$ of the chosen sides over the auxiliary sides (\cref{prop:signed-sides}(ii)), $s\in\{{\rm L},{\rm R}\}$.
 Further $T_{QE}=S_E+A_E$ and
 \[
  T=S_E+A_E+S_{\rm oth}+A_{\rm oth}+T_K+S_{\rm u}(a)+A_{\rm u}(a)+10^{-12}.
 \]
\item \emph{Invariant price.} With $\pi_i=(f_1+h_1)_i$ the polarization coefficient (an upper end) of \cref{sec:envelopes} for the piece,
 \[
  \Uinv=\sum_i\Bigl[2v_{{\rm D},i}I^{\rm ex}_{2,i}
  +\frac{12w_{{\rm S},i}}{1-3/X_i}I^{\rm ex}_{3,i}\Bigr],
 \]
 the price at $a_0$ assigned to the invariant tags $E_{\rm add}+H$ of the
 bracket (\cref{lem:oct24:invariant-background}(v)); it is an exact rational,
 $\Uinv=0.84688441321358980128331130\ldots$, and with the price
 $\sum_i2v_{{\rm D},i}\pi_iI^{\rm ex}_{3,i}$ assigned to the generic tag $G$ it
 satisfies $\Uinv+\sum_i2v_{{\rm D},i}\pi_iI^{\rm ex}_{3,i}\le U_0'$. The
 \emph{background factor} is $e^{t\Uinv}$; the tag $G$ is not invariant and is
 paid in $U_0$ only.
\item \emph{Leading class costs.} For each of the sixteen envelope classes of
 \cref{cor:envelope-costs} (the classes $(n,m)$ with $n\ge1$,
 $n+m\in\{2,4,6\}$; the aligned and the crossed matching pairs, written
 $(EE,{\rm al})$ and $(EE,{\rm cr})$; and $(0,4)$, $(0,6)$), with the caps
 $\kappa_H\ge\kappa_b\ge0$ (generic and invariant backgrounds, \cref{cor:envelope-costs}), the coupling mass $M_{n,m}$ and the
 invariant orbit mass $M^{\rm inv}_{n,m}$ of its envelope
 (\cref{tab:data:envelopes}),
 \[
  C_{n,m}=\kappa_He^{U_0}-(\kappa_H-\kappa_b)e^{t\Uinv}+M_{n,m}R_{\rm err}
  -M^{\rm inv}_{n,m}e^{t\Uinv}t\delta_{\rm s},
 \]
 where
 $\delta_{\rm s}=\frac12D'-\sum_i\operatorname{up}\bigl((v_{{\rm D},i}+v_{{\rm S},i})
 s_{\rm fold}(X_i)I^{\rm hi}_{4,i}\bigr)$ is the price saving at $a_0$ of one
 folded tag, its sharp price minus its symmetric price (\cref{lem:orbit-coupon}),
 $s_{\rm fold}$ being the symmetric coefficient of
 \cref{lem:record:one-stage-folded-e}.
\item \emph{Savings.} With the remainder polynomial
 $\mathcal R_{n,m}=(A^{\rm lead}_F+A^{\rm rem}_F)^n(A^{\rm lead}_E+A^{\rm rem}_E)^m
 -(A^{\rm lead}_F)^n(A^{\rm lead}_E)^m$, the \emph{saving} of a class is
 \[
  D_{n,m}=\frac{BA_F^nA_E^m-C_{n,m}(A^{\rm arm}_F)^n(A^{\rm arm}_E)^m
  -B\mathcal R_{n,m}}{n!\,m!}
 \]
 for the twelve classes with $n\ge1$, $n+m\in\{2,4,6\}$ of (7) and for $(0,4)$, $(0,6)$; for the two
 matching pairs, $D_{EE,{\rm al}}$ and $D_{EE,{\rm cr}}$ are the same expression
 with $(n,m)=(0,2)$, the cost of the pair class and the divisor $4$ in place of $2$, since the aligned and crossed pairs each have total $A_E^2/4$ (proof of \cref{prop:closing-data}(d)). A class is replaced only if the lower end of its interval enclosure (\cref{cert:closing}) is positive;
 otherwise its saving is omitted. All sixteen savings are positive.
\item \emph{Variance credit.} $D_{\rm var}$ is \eqref{eq:variance:dvar} of \cref{cor:variance:final} evaluated with the following substitutions: the rows of (3) for its record weights, multiplied by the repricing factor $t$; the restoration function evaluated at the piece caps of (1) with the error $R_{\rm err}$ of (2); and the gains given by \cref{prop:background-variance} and \cref{lem:pool-coupons} (including their terms pooled over the pieces from $2^{31}$). Its skew pair credits are those of the fold records, with the cap
 $\kappa_{FF}$, and of the aligned matching pairs, with the invariant cap of
 their envelope and the gains proved only for the selected laws, and its credits of degrees three
 and four include those of the pure quartic classes $(4,0)$ and $(0,4)$ on their raw-arm prices
 (\cref{lem:class-credit}).
\item \emph{Discount and reserved cost.}
 \[
  D=D_{\rm var}+\sum_{n\ge1}D_{n,m}+D_{EE,{\rm al}}+D_{EE,{\rm cr}}+D_{0,4}+D_{0,6},
  \qquad
  \mathcal E=B\bigl(e^S\cosh A-1\bigr)-D.
 \]
\item \emph{Mean credit.} $m=1+m_F+m_Q-R_{\rm err}e^{T_K+T_{QE}}$, with the
 coupons $m_F$ (empty matching tag set) and $m_Q$ (nonempty matching tag set)
 of \cref{prop:record:mean-credit} at the piece caps, the late-pool coupons of
 degrees five and six included (\cref{lem:pool-coupons}).
\item \emph{Index terms.} $\Kint=\PushK$. For the conductor moment use the certified upper bound
 \[
 R_I\le R_I^+=1.247139907406884705556794769493
 \]
 for the finite depth-$64$ moment $M_{\rm pad}(\Kint)$ (every later occurrence of $R_I$ in an upper bound means $R_I^+$), certified using
 \cref{prop:imp:padic-price,cor:imp:padic-majorant}
 (\cref{cert:index-moment}); and
 \[
  \begin{aligned}
  J&=Be^T\Bigl(\frac{R_I}2\Bigr)^6,\\
  W&=Be^{S+A}\Bigl[\Bigl(1+\frac{R_I}2\Kint^{-35/36}\Bigr)^6-1\Bigr].
  \end{aligned}
 \]
\item \emph{Closing number.}
 $\Theta=\frac{63}{64}(L_0+\mathcal E)+J+W+\frac{1-m}{64}$.
\end{enumerate}
\end{definition}

The prices of (3) are rounded per piece on $10^{-30}\Z$ and are sums of upper
prices of the records of the piece, by \cref{lem:middle-comparison-prices} and
the block charges; by \cref{lem:record:record-expansion}(iii) they may be
distributed over the blocks of the piece in proportion to the actual prices,
and the complete polynomial, which has nonnegative coefficients, is evaluated
at them. The endpoint coefficients are values at the piece start of admissible
majorants (\cref{lem:majorant}), so they bound the coefficients of every prime
and block of the piece. No upper estimate of a norm is used as a lower estimate
of that norm; the subtracted exponential $e^{t\Uinv}$ and the coupon
$M^{\rm inv}_{n,m}e^{t\Uinv}t\delta_{\rm s}$ are subtracted only inside
coefficientwise majorants (\cref{rem:completion-subtraction}).

The closing has the form of the comparison theorem; the work is in the
discount. The even skew part $B(\cosh A-1)$ of the reserved cost is far larger
than its symmetric part, and the envelope classes of \cref{sec:envelopes}
replace it, degree by degree, by class costs on raw-arm monomials: a class of
$n$ fold and $m$ matching records replaces its full-norm term $BA_F^nA_E^m$ by $C_{n,m}(A^{\rm arm}_F)^n(A^{\rm arm}_E)^m$ for its leading parts and by $B$ times the remainder polynomial $\mathcal R_{n,m}$ of (8) for the rest. The cost $C_{n,m}$ bounds the energy of the blocks outside the class: $\kappa_He^{U_0}$ for a generic outside background, lowered on the invariant
outside background, whose price $t\Uinv$ is assigned exactly, plus the
couplings of the envelope, each paying the error $R_{\rm err}$ of the bracket,
less the orbit coupon of \cref{lem:orbit-coupon}.

\begin{proposition}[The closing quantities bound the data of the comparison theorem]
\label{prop:closing-data}
For the construction of Part~V, at every scale $\Lsc\ge1$:
\begin{enumerate}[label=(\alph*)]
\item the totals $U_0$, $P$ and the tag totals of (CP2), and $R_{\rm err}$,
 $L_0$, $B$ of (CP2) and (CP3), are at most the numbers of
 \cref{def:closing-quantities}(2) (the tag totals at most $d$ and $v$);
\item the summed symmetric and skew prices of (CP4) are at most $S$ and $A$;
\item the summed price of the matching-envelope diagonals of (CP6) is at
 most $T$;
\item the complete polynomial of \cref{lem:record:record-expansion}, formed
 with the class bounds of \cref{cor:envelope-costs} and the credits of
 \cref{def:closing-quantities}(9), at the upper prices of (3)--(4), is at most
 $\mathcal E$;
\item the mean energies satisfy (CP7) with the $m$ of (11);
\item $R_I$ of (CP0)(d) is at most the value of (12).
\end{enumerate}
The following gates hold: every retained saving of (8) is positive; every
orbit coupon of (7) is at most $M_{n,m}R_{\rm err}$; every restoration gain lies
in $[0,1)$, and every restoration coefficient $\ell$ (the $l$ of \cref{lem:record:restoration}) at cap $C_X$ satisfies $0\le\ell\le C_X/2$; in
every piece $k^{\rm aux}_{{\rm L},i},k^{\rm aux}_{{\rm R},i}\ge0$ and the
protection gate $v_{{\rm D},i}\ge2w_{{\rm D},i}k_{{\rm D},i}$ (\eqref{eq:protection})
hold, and so do the matching gates of \cref{sec:matching-mean}; $D'\le V'$,
$d<6$ and $dv<36$; $B>1$; every credit of (9) on a raw-arm
monomial of a replaced class is below the cost of that class
(\cref{lem:class-credit}); and $0<a\le9/10000$. Hence the closing number of
\cref{thm:comparison} for this construction is at most $\Theta$.
\end{proposition}

\begin{proof}
(a) The first-stage atoms are the middle blocks (\cref{def:record:frame}); they
have $e_p=1$ and $\log q_\Blk\le L$ (\cref{lem:record:local-domain}(ii)). By
\cref{lem:record:middle-prices} and \cref{cert:record:middle-price}, $U_0'$,
$P'$, $D'$ and $V'$ bound the totals $U_0$, $P$, $D$ and $V$ at $a_0$, and
each of the two tags carries half of $D$ and of $V$
(\cref{def:record:frame}(e)), so that $D'/2$ and $V'/2$ bound both tag totals
at $a_0$. \Cref{cor:first-stage-repricing} gives at $a$ the bounds $tU_0'$
and $tP'$ for $U_0$ and $P$, the bounds $d$ and $v$ for the sharp and the
variation total of each tag, and $L_0\le e^{tP'}\tagser$,
$R_{\rm err}\le e^{tU_0'}\tagser$, $B\le e^{tU_0'}(1+\tagser)$, where, by the
monotonicity of the series (\cref{rem:cp:tags}), $\tagser$ may be taken at
the common values $d$, $v$ of the two tags. Since $D'\le V'$, $d\le v$, and
with $d<6$ and $dv<36$ \cref{cor:two-cut-series} at $n=4$ bounds this series
by $\tagser_4^+(d,v)$.

(b) \emph{Middle.} By \cref{lem:middle-comparison-prices} the symmetric and
skew prices of the records of a middle block $\Blk$ with least prime $b$, and
their raw-arm, leading and remainder auxiliary prices, are at most
$q_\Blk^{2a}b^{-k/2}$ times the coefficients of that lemma, $k$ the degree of
the price; the coefficients are evaluated at the piece start and bound those of
every block of the piece (\cref{lem:record:local-domain}(iii)). Summed over the
blocks of the piece at $a_0$, $\sum_\Blk q_\Blk^{2a_0}b_\Blk^{-k/2}\le\bar I_{k,i}$,
so the rows of (3) bound the prices of the records of the piece at $a_0$.
Every middle block has $\log q_\Blk\le L$, so at $a$ these prices are at most
$t$ times their values at $a_0$ (\cref{lem:record:repricing}); summing over
the pieces gives $S_{\rm mid}$ and $A_{\rm mid}$. \emph{Upper bands and tail:}
\cref{prop:record:upper}(ii), every term of the band ledger repriced from
$a_{\rm u}$ at its own logarithmic product (at most $L_{\rm band}\le886$), and
\cref{lem:tail-prices}(iv). \emph{Prefix:} the deleted exception (the coincident-label part of the prefix comparison) is a diagonal at the reserved cut $\iota$ and is included in $S$ only, at the total price, with $\mcl=8192$ the number of digit pairs, $2^{513-\mcl+1024\mcl a}=2^{-1380939795479/6103515625}<8\cdot10^{-69}$
(\cref{lem:frontier:iidenergy}), inside the allowance $10^{-12}$
(\cref{cert:record:frame-data}(f)), which also holds a second copy of the
tail prices (\cref{rem:prefix-digits}). These are all the atoms
(\cref{def:construction}(a)).

(c) This is \cref{prop:record:matching}(ii) with the chosen sides of
\cref{prop:signed-sides}. In a middle block the envelope $M_{s,\Blk}-\Nbr_\Blk$
consists of the chosen side diagonals, of price at most $T_K$
(\cref{prop:signed-sides}(iii) and the factor $t$), the chosen matching
diagonal, of price at most $T_{QE}$ (by the matching gates of \cref{sec:matching-mean}),
and the norm envelopes of the other record changes,
of price at most $S_{\rm oth}+A_{\rm oth}$; outside the middle the envelopes are the
norm envelopes of the upper, tail and prefix changes, of price at most
$S_{\rm u}(a)+A_{\rm u}(a)+10^{-12}$. The prices of the fold records do not
enter $T$.

(d) We use \cref{lem:record:record-expansion} with $B'$ the number $B$ of (2),
which bounds the energy factor of (CP3) by (a), at the per-block upper prices
and auxiliary prices obtained from the totals of (3)--(4) by
\cref{lem:record:record-expansion}(iii), the two records of one kind in a
block receiving equal shares; the auxiliary prices of a record are its
raw-arm, leading and remainder prices (\cref{lem:middle-comparison-prices}).
Every selection of a class has at most six atoms.

\emph{Variance classes.} The classes of \cref{cor:variance:final}(i) and (iii)
whose selections contain a symmetric factor have the class costs $B-g$ with
gains $g\in[0,1)$. By \cref{lem:record:record-expansion}(iv) the series
$\Psi_0=B'(e^{S(x)}\cosh A(x)-1)-\mathcal D_{\rm sym}(x)$, in which
$\mathcal D_{\rm sym}$ carries these gains, the monomials with a repeated
block receiving the gain of their piece multiset, below $1<B'$, is a complete
polynomial for them; its coefficients on the monomials in skew variables alone
are those of $B'(e^{S(x)}\cosh A(x)-1)$.

\emph{Envelope classes.} The sixteen classes of \cref{cor:envelope-costs}
satisfy (CP5) with the class bounds
$C_{n,m}\pi^{\rm arm}_\omega+B(\pi^{\rm tot}_\omega-\pi^{\rm lead}_\omega)$
(\cref{thm:envelopes}(i)), and $C_{n,m}\ge0$: since $\Uinv\le U_0'$,
$e^{U_0}\ge e^{t\Uinv}$, so that
$\kappa_He^{U_0}-(\kappa_H-\kappa_b)e^{t\Uinv}\ge\kappa_be^{t\Uinv}$, and the orbit
coupon is at most $M_{n,m}R_{\rm err}$ (\cref{lem:orbit-coupon}); hence
$C_{n,m}\ge\kappa_be^{t\Uinv}\ge0$. By \cref{lem:record:record-expansion}(v), with
the kinds of the remark after it (the fold records and the merged matching
records for the general classes and for $(0,4)$, $(0,6)$; the two matching
records $E_0$, $E_1$ separately for the pairs, the aligned and the crossed
pairs each having the total $A_E^2/4$), replacing these classes subtracts
from $\Psi_0$ polynomials whose values at the upper prices are the savings
of (8); the classes are disjoint and none contains a selection of a
variance class of the previous step.

\emph{Credits.} Inside the classes $(2,0)$, $(EE,{\rm al})$, $(4,0)$ and
$(0,4)$ the credits of \cref{lem:class-credit} are subtracted on the raw-arm
monomials, as the last part of \cref{lem:record:record-expansion}(v) allows:
every credit $cg$ satisfies $0\le cg<C_{n,m}$ (\cref{lem:class-credit}), so the
credited class bounds have nonnegative coefficients. Indeed the
gains $g$ are below one and $c\le\kappa_be^{t\Uinv}\le
C_{n,m}$: in the first three classes $c=\kappa_b$, and in class
$(0,4)$ the aligned variance cap satisfies the certified
inequality $c<\kappa_be^{t\Uinv}$ (the two sides are $0.667$
and $1.12$ to three digits; the certified values are in
\cref{tab:data:envelopes,cert:envelopes}). By
\cref{cor:variance:final}(iv) the subtractions of $\mathcal D_{\rm sym}$ and of
these credits together lower the polynomial at the upper prices by at least
$D_{\rm var}$, every coefficient staying nonnegative.

\emph{Conclusion.} The resulting series $\Psi$ is a complete polynomial
(\cref{lem:record:record-expansion}(ii)). At the upper prices the totals of
the price and auxiliary variables are those of (3)--(4), so
$\Psi(\hat\pi,\hat y)\le B(e^S\cosh A-1)-(D-D_{\rm var})-D_{\rm var}=\mathcal E$.
This proves the complete bound; it does not assert that an isolated formal
discount bounds a distinct-block discount from below
(\cref{rem:completion-subtraction}).

(e) \Cref{prop:record:mean-credit} at the piece caps, with the error exponent
$T_K+T_{QE}$; its gates are among those listed. (f)
\Cref{prop:imp:padic-price,cor:imp:padic-majorant} and
\cref{cert:index-moment}.

The gates are asserted by the author assembly of \cref{cert:closing}: the
orbit coupons by \path{research/advance_2026_10_24/global/closing_beta.py},
lines 204--205; the savings by lines 223, 228, 234 and 240 there (a general or
quartic class is retained only if its saving is positive, and the classes
$(2,0)$ and the matching pairs must be); the series and $D'\le V'$, $d<6$,
$dv<36$ by \path{research/advance_2026_10_24/global/grouped_first_stage.py},
lines 27--37, which also checks that $\tagser_4^+$ does not exceed the
one-sharp bound of \cref{rem:cp:tags}; the regime $0\le\ell\le C_X/2$ by
\path{research/continuation_2026_10_24/envelopes/restoration.py}, lines 36--37;
the gains by lines 148 and 160 of \texttt{closing\_beta.py} and line 31 of
\path{research/advance_2026_10_24/envelopes/beta_higher_variance.py}; the caps
$1\le C_i\le C$ by line 16 of
\path{research/advance_2026_10_24/envelopes/piece_cap_restoration.py}; and the
sides $k^{\rm aux}_{s,i}\ge0$ by line 42 of \texttt{closing\_beta.py} and line
129 of \path{research/advance_2026_10_24/global/piece_data_beta.py}; and the
matching gates by line 112 of \texttt{closing\_beta.py}. The
protection gate is asserted on all $33$ selected native pieces by
\path{research/improvements_2026_10_25/envelopes/check_gap_native.py}, rerun
on the rebuilt source of \path{research/push_2026_10_25/design.json};
$0<a\le9/10000$ is asserted by the selected design verifier.
Moreover, $B>1$ since $U_0>0$, and the credit gates follow from the
bound $C_{n,m}\ge\kappa_be^{t\Uinv}$ shown in (d).

Finally, $\Theta$ is the closing number of \cref{thm:comparison} at $k=6$,
$d=72$, $\theta=1$ for $\mathcal E\ge\Psi(\hat\pi,\hat y)$, with $J$ and $W$ as in
\cref{prop:oct24-index-cost}. By \cref{rem:cp:monotone} the closing number
is nondecreasing in $L_0$, $B$, $S$, $A$, $T$, $R_I$, the prices and the
auxiliary prices, and nonincreasing in $m$; so the bounds (a)--(f) show that the
closing number of the construction is at most $\Theta$.
\end{proof}

\Cref{tab:closing-values} lists the values. The upper bands and tail
contribute $S_{\rm u}(a)\approx0.0138662$ and
$A_{\rm u}(a)\approx0.00668235$. The class savings remove about $0.456895$
and the variance credits about $0.00995856$ from the reserved cost. The
reserved contribution $\frac{63}{64}(L_0+\mathcal E)+\frac{1-m}{64}$ is
about $0.748200$, the matching term $J$ about $0.251761$, and the wrap $W$
about $3.8345\cdot10^{-5}$. These rounded descriptions are explanatory only.

\begin{table}[ht]
\centering
\caption{The closing quantities, rounded outward from the certificate's
enclosures ($<$ upper bound, $>$ lower bound, $\ldots$ exactly defined).
The certificate evaluates $\Theta$ from the full enclosures, not from these
values.}
\label{tab:closing-values}
\begin{tabular}{@{}llll@{}}
\toprule
quantity & value & quantity & value\\
\midrule
$t$ & $1.0000000616512\ldots$ & $D-D_{\rm var}$ & $>0.456895419675$\\
$C$ & $<2.753348440327$ & $D_{\rm var}$ & $>0.009958560315$\\
$U_0$ & $<0.846908245753$ & $D$ & $>0.466853979991$\\
$B$ & $<2.332628665480$ & $\mathcal E$ & $<0.768964199811$\\
$R_{\rm err}$ & $<0.000204255886$ & $m_F$ & $>0.400172842378$\\
$L_0$ & $<0.000476332373$ & $m_Q$ & $>0.190134971388$\\
$S_{\rm u}(a)$ & $<0.013866153512$ & $m$ & $>1.589940741988$\\
$A_{\rm u}(a)$ & $<0.006682351376$ & $R_I$ & $<1.247139907407$\\
$S$ & $<0.014554805967$ & $J$ & $<0.251761455211$\\
$A$ & $<0.969271826335$ & $W$ & $<0.000038344989$\\
$T$ & $<0.607496889411$ & $\Theta$ & $<\PushTheta$\\
\bottomrule
\end{tabular}
\end{table}

\subsection{Outward-interval evaluation: \texorpdfstring{$\Theta<\PushTheta$}{Theta<\PushTheta}}\label{subsec:closing-certificate}

It remains to evaluate the explicit number $\Theta$ defined in
\cref{subsec:closing-quantities}. This is a finite computation:
every price is a sum of explicit terms at finitely many prime
thresholds, evaluated in outward-interval arithmetic so that
rounding can only enlarge the result. The certificate states the
outcome and identifies the programs and data files that produce
it; \cref{app:certificates} explains what such a check does and
does not prove.

\begin{certificate}[The closing]\label{cert:closing}
The selected rational design
\path{research/push_2026_10_25/design.json} has
\[
 a=\PushA,\qquad \eps=a/71=\PushExponent,\qquad \Kint=\PushK,
\]
and the outward evaluation of \cref{def:closing-quantities} obeys
$\Theta<\PushTheta<1$. The entry point
\path{research/push_2026_10_25/verify.py} invokes the selected adapter
\path{research/improvements_2026_10_25/global/selected.py} with this design.
Two independently written assemblies of the scalar totals use the
separate evaluators of \path{research/advance_2026_10_24/}, with the
invariant-background and root-restoration adapters of the later
checkpoints. The finite conductor law is supplied by
\path{research/improvements_2026_10_25/global/padic_moment.py}; its
independent check and the recomputation of the closing from it are
part of the selected verification. The outputs equal the stored
\path{research/push_2026_10_25/artifacts/author.json} and
\path{research/push_2026_10_25/artifacts/independent.json}; the enclosures of all $45$ compared scalar components intersect. The independent enclosure of $\Theta$ is
\begingroup\footnotesize
\begin{gather*}
 [0.99999999997325941882402031814341317621743080026301922476436792225715828504384175,\\
 0.99999999997325941882402031814341317621743080026301922476436792225715828504384968],
\end{gather*}
\endgroup
and the author's upper endpoint is below
$0.999999999973259418824020318144$.
Both are strictly below $\PushTheta$, giving a certified margin greater than
$2.673\cdot10^{-11}$. The archival launcher \path{paper/verify.py} restores
and authenticates the original source layouts before replaying their
certificates (\cref{app:certificates}). These finite checks do not replace
the analytic reductions or external theorems, whose proofs and assumptions
are given in the text.
\end{certificate}

\paragraph{How the certificate evaluates $\Theta$.}
The certificate evaluates every item of \cref{def:closing-quantities} from the sources whose hashes \cref{app:certificates} checks: the selected middle source
\path{research/push_2026_10_25/global/refined_L570_source.json}, its exact
charge evidence and native parent; the unchanged upper source
\path{research/further_2026_10_25/upper/upper_dense_refined.json}; and the
envelope bundle
\path{research/further_2026_10_25/envelopes/exact_background_pairs.json}.
The exact charges are regenerated at $a_0=\PushAnchor$, then repriced by
$t=e^{2(a-a_0)L}$; the proof is blockwise monotonicity of
$q_\Blk^{2a}$, as in \cref{lem:record:repricing}.
The root-restoration adapter supplies the actual marginal data of all seven tilted pieces of \cref{def:positive-root}, and the conductor adapter supplies the finite-depth upper
majorant. All first-stage totals, comparison prices, discounts and coupons
are recomputed from these sources. The formulas and their checks are listed
in \cref{subsec:cert:evaluation}.

\subsection{Proof of the baseline (\texorpdfstring{$\theta=1$}{theta=1}) by Theorem~C}\label{subsec:proof-main}

Theorem~C now applies directly: it only remains to point, for each
hypothesis (CP0)--(CP8), to the statement of Part~V that verifies
it for the construction, and to $\Theta<1$ from the preceding
subsection.

\begin{proof}[Proof of \cref{thm:baseline}]
By \cref{thm:comparison} it suffices to verify, for the construction of
Part~V with $k=6$ queries on $d=72$ slots and $\theta=1$, the hypotheses
(CP0)--(CP8) at every scale, the boundedness of the totals independently of
the scale, and $\Theta<1$.

\emph{Hypotheses.} \Cref{tab:hypotheses} lists for each hypothesis the
statement of Part~V that verifies it: (CP0) the slot array, the queries,
the atoms and the analytic stage, with
$a=\PushA\le9/10000<1/825=\sigma$,
$\sigma-a=65195254/201416015625$, and the gates of the stage on the whole
domain $0<a\le9/10000$
(\cref{cert:record:analytic-stage,cert:record:frame-data}), and the index law
(\cref{prop:oct24-index-cost}); (CP1) the prefix, middle, upper and tail
references, with \ref{it:record:R4} (\cref{lem:record:tail-mixing}); (CP2),
(CP3) the first stage and the law dominating it in a $2\times2$ block;
(CP4) the comparison law in every range, with \ref{it:record:R5}, the signed
fold records being finite signed combinations of nonnegative accepting
product laws, so that the quantities $n_\nu$ and $\mathfrak v_\nu$ of \cref{def:cp4} are finite; (CP5) the
class bounds of \cref{sec:envelopes} and the variance classes of
\cref{sec:variance}, whose selections have at most six blocks;
(CP6), (CP7) the matching
envelopes and the mean credit; (CP8) admissibility (\cref{def:admissible}),
robustness (\cref{lem:record:rec}) and summability (\cref{lem:tail-prices}(iv)).

\emph{Bounded totals.} Every total is a sum of certified per-atom upper
prices over finitely many atoms below $\Pfour$, independent of $N$, and over
the tail primes, whose prices are summable (\cref{lem:tail-prices}(iv)). So
\cref{prop:closing-data} holds at every scale with the same constants.

\emph{Closing inequality.} By \cref{prop:closing-data} the closing number of
\cref{thm:comparison} is at most $\Theta$, and $\Theta<\PushTheta<1$ by
\cref{cert:closing}. Robustness, the clause of (CP8), follows from
\cref{lem:record:rec}(i) as in \cref{rem:cp:robust-final}: every gate on a
quantity that the inflation changes is strict (the savings, the gains, the
gates $d<6$, $dv<36$, and $\Theta<1$), except the orbit-coupon gate, which is a
nonnegativity condition on a coefficient that only grows under inflation; the
restoration regime $0\le\ell\le C_X/2$ is kept by the inflation, which lowers $\ell$ and raises the effective cap to $(1+\delta)C_X$ (\cref{rem:cp:robust-final}(c)); and the side gates $k^{\rm aux}_{s,i}\ge0$, the
protection gate and $0<a\le9/10000$ concern fixed numbers of the construction.
So the function $\Theta(\delta)$ of \cref{lem:record:rec}
($\delta$ the inflation parameter, with $\Theta(0)=\Theta$) is
continuous at $\delta=0$, with value below one there.

\emph{Conclusion.} \Cref{thm:comparison} gives, for every $\Lsc\ge1$,
$\Yf(\Lsc,\cdot)\le C\Lsc^{-a}\alpha$ at every square-difference-free subset
of $[\Kint\Lsc]$ of density $\alpha$, hence
$s(N)\le C'N^{1-a/71}$ for every $N\ge1$ with an absolute constant $C'$. As
$a/71=\PushExponent=\eps$, this is \cref{thm:baseline}.
\end{proof}

\paragraph{Consequences and what to read next.}
\Cref{thm:comparison,rem:cp:monotone,tab:hypotheses};
\cref{lem:record:repricing,cor:first-stage-repricing,cor:two-cut-series,lem:record:one-stage-folded-e}
and \cref{cert:record:frame-data,cert:record:middle-price};
\cref{lem:record:middle-prices,lem:piece-caps,lem:middle-comparison-prices,prop:signed-bounds,lem:oct24:matching-star,lem:record:exact-marginal-e};
\cref{prop:record:upper,lem:tail-prices,lem:frontier:iidenergy};
\cref{prop:record:matching,prop:signed-sides};
\cref{lem:record:record-expansion,rem:completion-subtraction};
\cref{cor:envelope-costs,thm:envelopes,lem:class-credit,lem:orbit-coupon,lem:oct24:invariant-background};
\cref{cor:variance:final,prop:background-variance,lem:pool-coupons};
\cref{prop:record:mean-credit};
\cref{prop:oct24-index-cost,lem:index-moments,cert:index-moment};
\cref{cert:record:analytic-stage}; and \cref{lem:record:rec,rem:cp:robust-final}.

\paragraph{Hypotheses verified in this section.}
The row ``$\Theta<1$'' of \cref{tab:hypotheses}, $\Theta<\PushTheta$
(\cref{prop:closing-data,cert:closing}), and \cref{thm:baseline}.
